\documentclass[11pt]{article}

\usepackage[margin=1in]{geometry}
\usepackage{booktabs}
\usepackage{enumitem}
\usepackage[hidelinks]{hyperref}

\newcommand{\code}[1]{\texttt{#1}}

\title{Gabriel Rust Rewrite: Progress}
\author{Status report}
\date{\today}

\begin{document}
\maketitle

\section{Summary}

The rewrite is roughly \textbf{45\% done} by effort, measured against the
implementation plan in the design note (\code{docs/dag-tokens.tex}). This is a
rough estimate. The pieces built so far are the ones that were hardest to get
right: the token rules, duplicate protection, and alignment at a fan-out. What
remains is mostly plumbing (processes, transports, serialization, and the C
ABI), which takes more code but involves fewer subtle decisions.

The first useful milestone, a flow running end to end through local engines
(plan steps 1--3), is \textbf{reached}: a producer's frames go through the
client and the local server to a chain of local engines, and the engines'
outcomes return the producer's tokens. A flow is added from Rust with a
\code{FlowSpec}, or from C with the same spec in TOML.

Plan step 4 is done for the local server: the producer's app reads what
happened to each of its inputs, and a client that loses its local server
reports \code{Disconnected} and installs its flows again. The plan for
logging, metrics, and benchmarks is in \code{docs/evaluation.tex}. Logging is
done, every event now carries timestamps, and a benchmark binary runs a flow
with synthetic engines and records them. A first sweep of token counts
showed how to set them: with too few, engines idle or pending inputs freeze;
one spare token keeps every pending slot fresh (\code{docs/evaluation.tex}).

All 198 tests pass.

\section{Done}

\begin{description}[style=nextline, leftmargin=1.5em]
  \item[Producer (\code{producer.rs}, \code{envelope.rs})] Publishes inputs
        over iceoryx2 into loaned shared memory, and notifies the client. Each
        input carries an id chosen by the producer: a random session per
        producer start and a sequence number per input. Publishing returns
        the id, so the app can match events to its frames.
  \item[Client thread (\code{client.rs}, \code{client/worker.rs},
        \code{client/producers.rs})] Subscribes to producers and keeps the
        newest input of each. Runs one flow router per flow: applies the local
        server's events, admits inputs, and sends them on their crossings.
        Wakes on producers and on the server's events. Adding a flow needs its
        producers added, and a producer cannot be removed while a flow uses
        it; a dropped client removes its flows from the server. A frame
        reaches its engine 0.2\,ms after it is published at the median.
        Publishes every event about a producer's inputs for the producer's
        app, and tells the server about offers it skipped on crossings. Checks
        every second that the local server runs, and looks for gaps in its
        event stream; if either fails, it ends every offer there with
        \code{Disconnected} and installs its flows again. 16 end-to-end tests
        with a real server and engines, among them one that checks under load
        that every offer ends exactly once, and one that restarts the server;
        plus 4 for the producers' ports.
  \item[Flow spec and compiler (\code{spec.rs}, \code{client/compile.rs})]
        \code{FlowSpec}, the Rust description of a flow: producers with token
        counts, engines with their location, and edges with timeouts. The
        compiler checks it (names, a DAG, edges end at engines) and splits it:
        edges between engines go to the local server, and each producer's
        edges form one crossing. It derives the crossing and ticket
        deadlines from the timeouts. Remote engines are rejected until QUIC
        exists. From C, the spec is TOML: an edge list with the same shape as
        \code{FlowSpec}, with timeouts such as \code{"300ms"}. 13 tests.
  \item[Subscription (\code{subscription.rs}, \code{ffi.rs})] How a
        producer's app reads the events about its inputs, from
        \nolinkurl{gabriel/producer/<name>/result}: \code{wait} without
        polling, and \code{poll}, which reports gaps in the numbers as
        \code{Lost}. The app sees \code{Sent}, \code{Result} (with the result
        bytes), \code{Skipped}, \code{TimedOut}, and \code{Disconnected}, per
        input and engine. In C: \code{gabriel\_subscribe\_producer},
        \code{gabriel\_subscription\_wait}, \code{gabriel\_subscription\_poll}.
        5 tests.
  \item[Local client transport (\code{client/local\_server.rs},
        \code{server/local\_clients.rs}, \code{wire.rs})] Control calls
        (\code{InstallFlow}, \code{RemoveFlow}) over iceoryx2 request-response,
        encoded with postcard and carrying a protocol version; inputs on
        crossings over a server input service, which wakes the server; and one
        numbered event stream per client, which wakes the client. The server
        records which client installed each flow, and only that client may
        remove it. The input service also carries skip notices, which the
        server publishes as \code{Skipped} on the target engines' event
        services. 8 tests, plus the end-to-end tests above.
  \item[Edges (\code{routing/edge.rs})] \code{EdgeState}: an edge's capacity,
        inputs in flight with their deadlines, and the pending slot, where a
        newer input replaces an older one. No I/O; the owner collects actions
        (send, skipped, timed out). An edge can be suspended while its server
        cannot be reached. 13 tests.
  \item[Tokens (\code{routing/tokens.rs})] \code{TokenPool}: end-to-end tokens
        per producer, and a ticket per admitted input that keeps the input
        until it is done. Returns each token exactly once, even with fan-out,
        duplicates, and timeouts, and tells late results from duplicates.
        Checks its invariant in debug builds. 18 tests.
  \item[Flow router (\code{routing/flow.rs})] \code{FlowRouter}: setup of
        producers and client-owned edges; admission (an input is admitted when
        its producer has a token and one of its edges can send it right away);
        applying the end of an offer (early discard, forwarding after a
        \code{Pass} across an uplink, duplicate and late events, fan-out
        alignment, no forwarding after the ticket deadline); crossings, which
        keep their capacity after \code{Accepted} until their first target
        finishes or they time out; expiry of offers and tickets; reports of
        offers skipped or timed out on the client's edges; and lost
        connections, which end every offer on the server with
        \code{Disconnected} and stop the crossings until the client
        reconnects. 22 tests.
  \item[Local services (\code{ipc/types.rs}, \code{services.rs},
        \code{naming.rs})] The four iceoryx2 services of one machine, with
        fixed-size records: the engine's input (\code{InputHeader}), the
        engine's outcome (\code{OutcomeRecord}), and the engine's events
        (\code{EventRecord}, with one constructor per kind of event, and the
        flow it belongs to). The engine's input and outcome services never
        drop a sample, and have a small fixed buffer that does not depend on
        how many clients use the engine. Round trips through each service are
        tested.
  \item[Engine (\code{engine.rs}, \code{ffi.rs})] The engine side of the C
        ABI: create and destroy, \code{wait} (for inputs, without polling),
        \code{poll\_input}, and the outcomes \code{stop}, \code{pass}, and
        \code{emit} (with \code{new\_output}). Each outcome wakes the server.
        One input at a time. An input is one buffer with a payload type;
        Gabriel does not look inside it, so an engine that emits several tiles
        frames them itself. The engine names the type of its results when it
        is created, and of the new input when it emits. 14 tests.
  \item[Server router (\code{routing/server\_router.rs})]
        \code{ServerRouter}: installing and removing flows (checked against
        the registered engines, each with its own session); inputs arriving on
        crossings, offered to one edge per target engine; outcomes, which free
        the engine's place, end the offer, and forward \code{Pass} and
        \code{Emit} on the server's own edges, reporting the \code{Result}
        before its children's events; late, duplicate and misdirected
        outcomes; expiry; skip notices from clients. Engines pull their
        inputs: an engine is given one input at a time, and its edges keep
        their newest input in their pending slots and take turns. Lost inputs
        and dead engines free their places, so an engine's places cannot
        leak. No I/O. 30 tests.
  \item[Server dispatcher (\code{server/dispatch.rs})] Connects the server
        router to the engines' iceoryx2 services: delivers inputs and wakes
        the engine, waits for outcomes, reads them (checking their lengths),
        numbers and publishes each engine's events, and queues events for
        clients. A delivery that fails or reaches no engine, and a malformed
        outcome, lose their input in the router. Notices engines whose
        process stopped, by their iceoryx2 node, and cleans up after dead
        ones. 13 tests, with real engines.
  \item[Server thread (\code{server.rs}, \code{server/worker.rs})] Runs the
        dispatcher and the local client transport on its own thread, woken by
        engines' outcomes and clients' inputs, with commands from the
        \code{Server} handle over a channel. Checks for stopped engines every
        second. From an outcome to the engine's next
        input takes 46\,\textmu s at the median and 0.4\,ms at the 99th
        percentile. The C ABI creates a local server and adds and removes
        engines. 3 tests here; the end-to-end tests are the client's.
  \item[Services (\code{services.rs}, \code{naming.rs})] Besides the
        services above: an event per engine in each direction, so neither
        side polls; the server's control, input, and input-ready services,
        and each client's event stream and events-ready event; node names
        that identify engine, server, and client processes; and a
        backpressure handler that makes a send to a full buffer fail instead
        of hang.
  \item[Logging (\code{logging.rs})] With \code{tracing}: warnings for every
        failure the server and client threads carry on after, lost engines,
        and lost connections; info for engines and flows; each thread's
        events carry its name. Nothing per input. C apps print the log with
        \code{gabriel\_log\_init}. 1 test.
  \item[Timestamps (\code{clock.rs})] Every event says when it happened,
        and an event that ends a server's offer says when the offer was
        made. Both are read from \code{CLOCK\_MONOTONIC}, which every process
        on a machine shares, through \code{rustix}. The subscription and the C
        event view expose them. 3 tests, one of them end to end.
  \item[Benchmarks (\code{bench/})] \code{gabriel-bench} runs a flow
        from a TOML file with synthetic producers and engines on threads, and
        writes one CSV row per input and per event, a copy of the run's file,
        and the git commit. \code{analyze.py} reports, in queueing terms,
        load, utilization, waiting, idle time, tokens in use, and the age of
        information. Three example runs. 4 tests.
  \item[Design note (\code{docs/dag-tokens.tex})] Covers the architecture,
        ownership of edges, installing a flow, flow control, forwarding,
        copies, result events, token accounting, worked examples (early
        discard, multi-tier fusion, feedback), and disconnected operation.
\end{description}

\section{Remaining work}

\begin{center}
\small
\begin{tabular}{@{}rp{0.5\linewidth}p{0.36\linewidth}@{}}
\toprule
Step & What & Status \\
\midrule
1  & Routing logic: edges, tokens, flow router
   & \textbf{done} \\
2  & Server with engines over iceoryx2, and the new engine ABI
   & \textbf{done} for a local server; the rest of the engine ABI is
     step 7 \\
3  & Client--local server protocol over iceoryx2 (\code{InstallFlow}, inputs)
   & \textbf{done} \\
4  & Result events back to producers
   & \textbf{done} for the local server \\
5  & Wire format: postcard messages
   & done for control calls \\
6  & The same protocol over QUIC to the remote server
   & \textbf{done}, tested over loopback: flows split between the servers,
     uplinks, remote crossings with a configured capacity; the bench runs
     it over netem on loopback; an \code{Emit} does not cross yet \\
7  & Result edges, engine inboxes, \code{Derived} results (fusion, feedback)
   & not started \\
8  & Crossings with several targets
   & routing done on both sides; not yet carried over a transport \\
9  & Disconnected operation: priority producers, deduplication, stream
     priorities
   & not started \\
10 & Zero copy on one machine
   & optional \\
\bottomrule
\end{tabular}
\end{center}

\paragraph{Cross-cutting work} not tied to one step:
\begin{itemize}[nosep]
  \item the flow compiler, which splits a flow for the local server only so
        far;
  \item the C ABI for all new calls, generated with cbindgen;
  \item reconnecting after a lost connection;
  \item end-to-end tests across processes.
\end{itemize}

\section{Decisions since the last report}

\begin{itemize}[nosep]
  \item \textbf{Copies are fine for now.} A frame is copied at each service
        boundary; zero copy on one machine stays an optional last step. A
        first measurement shows the copies matter for raw frames: they are
        most of the 6.2\,ms round trip of a 4K frame
        (\code{docs/evaluation.tex}). Kept as they are until the remote
        benchmarks are done.
  \item \textbf{Outcomes go through the server.} Engines never publish events.
        The server frees the edge, creates the follow-on offers, and then
        publishes the event, so each event service has one publisher and one
        order. This is also what lets the same engine run on the cloudlet.
  \item \textbf{Gabriel does not know about parts.} An input is one buffer
        with a payload type. An engine that emits several tiles lays them out
        itself, and the next engine reads them by the type.
  \item \textbf{Payload types name how to read bytes.} They are set when an
        engine is created, and when it emits. Result events carry them.
  \item \textbf{Fixed-size records inside, a schema outside.} Records on
        iceoryx2 have a fixed layout. Apps are free to use protobuf for the
        bytes they send.
  \item \textbf{Every client edge into a server is a crossing}, including a
        producer's edge to a local engine. A server receives inputs from a
        client in one way, and offers each on its own edge per target engine.
  \item \textbf{An engine's input buffer is fixed.} iceoryx2 cannot resize a
        service, and clients come and go at runtime. So the server delivers
        a fixed number of inputs to an engine before their outcomes arrive.
  \item \textbf{Engines pull their inputs.} A server's edges have no
        capacity. Each holds its newest input in its pending slot until the
        engine has room, and the engine's edges take turns. Sending as soon
        as an edge had room would queue stale inputs at the server behind
        newer ones.
  \item \textbf{An engine gets one input at a time, and local crossings
        have capacity 3,} so an engine always starts on the newest input
        there is. A second input given ahead of time could not be replaced,
        and every slow engine would start on inputs one processing time old,
        doubling its latency. Capacity 3 (the delivery limit plus 2) is the
        smallest that keeps the server's pending slots updating when every
        target of a crossing is busy. The cost is a gap of tens of
        microseconds after each outcome, so the server thread must wake on
        the engines' notifications. Remote crossings face the same choice
        against bandwidth; it is an open question.
  \item \textbf{Deferred:} gating an edge on whether its engine's output can
        go anywhere, which saves compute before a slow engine; and pending
        slots that hold more than one input, for streams where every input
        matters. Both only extend when an edge may send, and what its
        pending slot holds.
  \item \textbf{Events name their flow}, because one input can enter several
        flows. A late result is not marked; a reader that saw its offer end
        knows.
  \item \textbf{An engine's death ends its delivered offers with
        \code{TimedOut} and its pending ones with \code{Skipped}.} The server
        thread notices it by the engine's iceoryx2 node: the node it saw last
        is gone, or none runs while inputs are delivered. It compares nodes
        rather than looking for dead ones, because a restarting engine cleans
        up dead nodes itself.
  \item \textbf{A full input or outcome buffer fails the send.} iceoryx2
        would otherwise make the sender wait, so a broken delivery bound
        would hang the server or an engine.
  \item \textbf{The server wakes on outcomes and on clients' inputs, and
        polls commands and control calls every 10\,ms.} Both are rare.
        iceoryx2 has no wake-up for a reply, so the client polls for it every
        millisecond.
  \item \textbf{Control calls are encoded with postcard}, the only messages
        without a fixed layout; inputs and events keep fixed records. Every
        call carries a protocol version, which the server checks.
  \item \textbf{Each local client has its own event stream}, so it reads only
        its own events, and a slow client fills only its own buffer. The
        stream is numbered, and drops its oldest events when full.
  \item \textbf{The spec gives timeouts, not deadlines or capacities.} Local
        crossings have capacity 3, and the compiler derives each crossing's
        and ticket's deadline from the engines' timeouts, as upper bounds.
  \item \textbf{Flow specs are TOML from C}, as an edge list with the same
        shape as \code{FlowSpec}, parsed by serde. Misspelled fields are
        errors.
  \item \textbf{A flow's inputs are the ones its producers publish on it}, by
        the flow named in each frame.
  \item \textbf{The producer's app sees events per input and engine, not per
        offer.} The client leaves out \code{Accepted}, reports a crossing
        offer that ends on the client once per target engine, and a lost
        connection once per flow. The result service has a ready event, so
        the app waits without polling.
  \item \textbf{A lost connection ends every offer on the server with
        \code{Disconnected}}, and the crossings to it stop until the flow is
        installed again under a new session. The offer pending in a crossing
        is skipped, so nothing old is sent after the server is back. The
        client checks the server's node every second, like the server checks
        its engines, and treats a gap in its event stream the same way.
  \item \textbf{Skip notices travel with inputs.} A crossing offer skipped on
        the client is sent to the server as a header with no payload, on the
        same service as inputs, and the server tells only the target engines'
        subscribers.
  \item \textbf{Logging uses \code{tracing}, and Gabriel installs no
        subscriber by itself.} Rust apps choose their own; C apps call
        \code{gabriel\_log\_init}. Warnings for failures, info for
        configuration changes, nothing per input.
  \item \textbf{Timestamps travel on the events}, instead of each process
        writing a file of them. The app already gets every event, so it sees
        an input's whole trip in one place, and a process killed during a
        benchmark loses nothing. The benchmark binary will write one CSV row
        per event; Parquet is not worth its large crates at our sizes
        (\code{docs/evaluation.tex}).
  \item \textbf{Tokens bound the flow, not freshness.} The token count must
        leave one token over when every engine and pending slot holds an
        input, or a pending input can never be replaced (398\,ms of waiting
        at p99, against 20\,ms with one more token). More tokens cost only copies of
        inputs that get skipped. A client-side handover of tokens was
        considered and rejected: it is not atomic, since the client learns
        late that an input was sent; a spare token gives the same result with
        the server deciding. The network is bounded by crossing capacities.
  \item \textbf{The objective is the age of information.} Gabriel minimizes
        the waiting it adds while keeping the slowest engine busy; AoI
        combines both, and every comparison reports waiting, idle time, and
        waste to explain it. Measures use queueing terms: load, utilization,
        Little's law, and the knee of the token count.
  \item \textbf{Tokens are a cap, not the control.} Below the knee they cost
        throughput, at exactly the slots they freeze them, and above that
        they are unused. Per-hop limits keep inputs fresh on the device; over
        the network, the crossing capacity is the control to tune. Tickets
        and token accounting stay as they are.
  \item \textbf{An input that will get no outcome is lost:} its offer ends
        with \code{TimedOut} at once and its place is freed. This covers a
        failed delivery, an engine that has not started, and a malformed
        outcome.
\end{itemize}

\section{Known gaps}

A review of the design against its goals and failure modes (``Known gaps'' in
the design note) found no gap that needs a new mechanism. Each one is fixed by
ending offers the way the design already does, or by detecting dead iceoryx2
nodes or gaps in sequence numbers.
\begin{itemize}[nosep]
  \item \textbf{A crashed local client's flows stay installed.} One that
        stops cleanly removes them. Fix: the server already records which
        client installed each flow, and would remove them when that client's
        node dies.
  \item \textbf{A skip notice or an event for the app can be lost} if a send
        fails. Nothing waits on it for tokens, so only a subscriber may wait
        longer than needed.
\end{itemize}

\section{Next steps}

\begin{enumerate}[nosep]
  \item \textbf{Compute the token count in the compiler}: engines plus
        server edges on the producer's paths, plus one (``Choosing $C$'' in
        the design note). Use it when the spec gives no count, and warn when
        a spec gives fewer.
  \item \textbf{Measure compute waste on the device} (later): engines that
        share a GPU, and whether gating on downstream readiness lowers the
        AoI (``Are tokens the right tool?'' in the design note). The
        synthetic engines only sleep, so an input that is later skipped
        costs nothing, and newest wins looks free. The bench needs an option
        in the run file that puts some engines behind one shared lock, in
        place of the GPU. Then measure the compute spent on inputs that are
        later skipped, and the AoI with and without gating.
  \item \textbf{Benchmarks} on the local server, as planned in
        \code{docs/evaluation.tex}: longer runs, sweeps of one knob at a
        time, and the baselines. Compare each input's age with the newest
        input that had reached its slot.
  \item \textbf{With the remote server}: sweep the crossing capacity by
        hand over slower links and with loss, then an adaptive crossing
        capacity, then predictive pull, compared on emulated links
        (\code{bench/runs/remote.toml} is the starting point).
  \item \textbf{The rest of the engine ABI}: \code{subscribe} to engines'
        event services (reusing \code{Subscription}), \code{publish\_result},
        and \code{wait} for subscribed events too (plan step 7).
\end{enumerate}

The first version copies a frame at each service boundary
(``Copies'' in the design note). Zero copy on one machine stays plan
step 10. The comparison with the original Gabriel showed that the copies
matter for raw frames: the server spends 2.5\,ms of CPU per 4K frame on two
of them. A smaller step comes first: hold received samples instead of copying
them into reference-counted buffers, which saves half the copies on one
machine without changing behavior.

\paragraph{Small open items.}
\begin{itemize}[nosep]
  \item A few error messages could name the producer or edge they are about.
  \item Nothing checks that an engine's payload type matches what its
        incoming edges carry. A mismatch shows up as garbage in the engine.
  \item The sizes of the fixed records (names, the offers per event, the
        engine input and outcome buffers) are guesses.
  \item Producers set no payload type yet, so inputs from producers reach
        engines with an empty one.
\end{itemize}

\end{document}
